% Part: computability
% Chapter: recursive-functions
% Section: halting-problem

\documentclass[../../../include/open-logic-section]{subfiles}

\begin{document}

\olfileid{cmp}{rec}{hlt}
\olsection{مشكلة التوقف}

\emph{مشكلة التوقف} هي، على العموم، أن نبت هل ينتهي حساب دالة عند
الدخل~${\OLId{latin-ordinary}{n}}$ إلى نتيجة أم لا، متى أعطينا مواصفتها~${\OLId{latin-ordinary}{e}}$،
كبرنامجها مثلًا، وأعطينا العدد~${\OLId{latin-ordinary}{n}}$. ومن نتائج آلان تورنغ
الشهيرة أن هذا السؤال العام لا تجيب عنه دالة قابلة للحساب؛
فالدالة
\[
{\OLId{latin-ordinary}{h}}({\OLId{latin-ordinary}{e}}, {\OLId{latin-ordinary}{n}}) =
\begin{cases}
{\OLMathNumeral{1}} & \text{إذا توقف الحساب ${\OLId{latin-ordinary}{e}}$ عند الدخل ${\OLId{latin-ordinary}{n}}$}\\
{\OLMathNumeral{0}} & \text{في غير ذلك،}
\end{cases}
\]
ليست قابلة للحساب.

وأما في الدوال العودية الجزئية، فالفهرس~${\OLId{latin-ordinary}{e}}$ الوارد في
مبرهنة كليني للصورة السوية يقوم مقام مواصفة البرنامج. فمتى كانت ${\OLId{latin-ordinary}{f}}$ دالة
عودية جزئية، عُدّ كل~${\OLId{latin-ordinary}{e}}$ يحقق معادلة تلك المبرهنة فهرسًا
لـ~${\OLId{latin-ordinary}{f}}$. ولكل عدد~${\OLId{latin-ordinary}{e}}$ تقضي مبرهنة الصورة السوية بأن
\[
\cfind{{\OLId{latin-ordinary}{e}}}({\OLId{latin-ordinary}{x}}) \simeq {\OLId{latin-looped}{U}}(\mu {\OLId{latin-ordinary}{s}} \; {\OLId{latin-looped}{T}}({\OLId{latin-ordinary}{e}}, {\OLId{latin-ordinary}{x}}, {\OLId{latin-ordinary}{s}}))
\]
دالة عودية جزئية، وأنه لكل دالة عودية جزئية
${\OLId{latin-ordinary}{f}}\colon\Nat\to\Nat$ يوجد ${\OLId{latin-ordinary}{e}}\in\Nat$ بحيث
$\cfind{{\OLId{latin-ordinary}{e}}}({\OLId{latin-ordinary}{x}})\simeq {\OLId{latin-ordinary}{f}}({\OLId{latin-ordinary}{x}})$ لكل~${\OLId{latin-ordinary}{x}}\in\Nat$. بل لكل
${\OLId{latin-ordinary}{f}}$ من هذا القبيل ما لا يتناهى من الفهارس~${\OLId{latin-ordinary}{e}}$. وتعريف
\emph{دالة التوقف}~${\OLId{latin-ordinary}{h}}$ هو
\[
{\OLId{latin-ordinary}{h}}({\OLId{latin-ordinary}{e}}, {\OLId{latin-ordinary}{x}}) =
\begin{cases}
{\OLMathNumeral{1}} & \text{إذا كانت $\cfind{{\OLId{latin-ordinary}{e}}}({\OLId{latin-ordinary}{x}}) \fdefined$}\\
{\OLMathNumeral{0}} & \text{في غير ذلك.}
\end{cases}
\]
وتكون ${\OLId{latin-ordinary}{h}}({\OLId{latin-ordinary}{e}},{\OLId{latin-ordinary}{x}})={\OLMathNumeral{0}}$ إذا وفقط إذا كانت
$\cfind{{\OLId{latin-ordinary}{e}}}({\OLId{latin-ordinary}{x}})\fundefined$.

% تصحيح مصدر (C226-01): كل عدد طبيعي فهرس في التعداد المعتمد، فحُذف الفرع المستحيل.
\textbf{تنبيه تصحيحي على الأصل.} يحصي الترميز المعتمد هنا دالة
عودية جزئية~$\cfind{{\OLId{latin-ordinary}{e}}}$ لكل عدد طبيعي~${\OLId{latin-ordinary}{e}}$؛ فحُذف الفرع الذي يفترض
وجود عدد طبيعي لا يكون فهرسًا أصلًا.

\begin{thm}
\ollabel{thm:halting-problem}
دالة التوقف~${\OLId{latin-ordinary}{h}}$ ليست عودية جزئية.
\end{thm}

\begin{proof}
لنفرض عودية ${\OLId{latin-ordinary}{h}}$ الجزئية؛ فيجوز حينئذ أن نعرّف
\[
{\OLId{latin-ordinary}{d}}({\OLId{latin-ordinary}{y}}) =
\begin{cases}
{\OLMathNumeral{1}} & \text{إذا كان ${\OLId{latin-ordinary}{h}}({\OLId{latin-ordinary}{y}}, {\OLId{latin-ordinary}{y}}) = {\OLMathNumeral{0}}$}\\
\umin{{\OLId{latin-ordinary}{x}}}{{\OLId{latin-ordinary}{x}} \neq {\OLId{latin-ordinary}{x}}} & \text{في غير ذلك.}
\end{cases}
\]
ولا عدد ${\OLId{latin-ordinary}{x}}$ يحقق ${\OLId{latin-ordinary}{x}}\neq {\OLId{latin-ordinary}{x}}$؛ ولذلك لا تتعين
$\umin{{\OLId{latin-ordinary}{x}}}{{\OLId{latin-ordinary}{x}}\neq {\OLId{latin-ordinary}{x}}}$، ومن ثم تكون ${\OLId{latin-ordinary}{d}}({\OLId{latin-ordinary}{y}})\fundefined$ إذا وفقط إذا كان
${\OLId{latin-ordinary}{h}}({\OLId{latin-ordinary}{y}},{\OLId{latin-ordinary}{y}})\neq{\OLMathNumeral{0}}$. فنحصل من التعريف على الأمرين الآتيين:
\begin{enumerate}
\item ${\OLId{latin-ordinary}{d}}({\OLId{latin-ordinary}{y}})\fdefined$ إذا وفقط إذا كان ${\OLId{latin-ordinary}{h}}({\OLId{latin-ordinary}{y}},{\OLId{latin-ordinary}{y}})={\OLMathNumeral{0}}$، أي إذا وفقط إذا
  كانت $\cfind{{\OLId{latin-ordinary}{y}}}({\OLId{latin-ordinary}{y}})\fundefined$.
\item ${\OLId{latin-ordinary}{d}}({\OLId{latin-ordinary}{y}})\fundefined$ إذا وفقط إذا كان ${\OLId{latin-ordinary}{h}}({\OLId{latin-ordinary}{y}},{\OLId{latin-ordinary}{y}})={\OLMathNumeral{1}}$، أي إذا وفقط إذا
  كانت $\cfind{{\OLId{latin-ordinary}{y}}}({\OLId{latin-ordinary}{y}})\fdefined$.
\end{enumerate}
ومن عودية ${\OLId{latin-ordinary}{h}}$ الجزئية تلزم عودية ${\OLId{latin-ordinary}{d}}$ الجزئية؛ فتجعل لها
مبرهنة كليني للصورة السوية فهرسًا~${\OLId{latin-ordinary}{e}}_{\OLId{latin-ordinary}{d}}$، أي
$\cfind{{\OLId{latin-ordinary}{e}}_{\OLId{latin-ordinary}{d}}}\simeq {\OLId{latin-ordinary}{d}}$. ولكن التعادلين السابقين يعطيان
\[
{\OLId{latin-ordinary}{d}}({\OLId{latin-ordinary}{e}}_{\OLId{latin-ordinary}{d}})\fdefined \quad\text{إذا وفقط إذا}\quad
\cfind{{\OLId{latin-ordinary}{e}}_{\OLId{latin-ordinary}{d}}}({\OLId{latin-ordinary}{e}}_{\OLId{latin-ordinary}{d}})\fundefined
\quad\text{إذا وفقط إذا}\quad {\OLId{latin-ordinary}{d}}({\OLId{latin-ordinary}{e}}_{\OLId{latin-ordinary}{d}})\fundefined\text{،}
\]
وهو محال. فبطل فرض عودية~${\OLId{latin-ordinary}{h}}$ الجزئية.
\end{proof}

\end{document}
